\chapter{Hilbert's Nullstellensatz}

THEOREM 1. --- Let A be an integral domain, K be its field of fractions, and L be a commutative K-algebra. Suppose that the A-algebra L is generated by finitely many elements and that L is a field.

\begin{enumerate}
\def\labelenumi{\alph{enumi})}
\item
  The degree of L over K is finite.
\item
  There exists a nonzero element \(a\) of \(\mathrm{A}\) such that \(\mathrm{K}\) is equal to \(\mathrm{A}[a^{-1}]\) .
\end{enumerate}

Let \(S\) be a generating subset of the A-algebra \(L\) . We use induction on the cardinal of \(S\) .

First, suppose that the elements of S are all algebraic over K. The degree of L over K is then finite (V, §3, No.~2, p.~18, Theorem 2). Let \((e_i)_{i\in I}\) be a basis of L over K. There exists a nonzero element \(a\) of A such that the coordinates with respect to this basis of the elements of S, of the element 1, and of the elements \(e_i e_j\) for \(i\) , \(j\) in I belong to \(\mathrm{A}[a^{-1}]\) . The set of linear combinations \(\sum_{i\in I}a_i e_i\) with \(a_i\in \mathrm{A}[a^{-1}]\) for every \(i\) is then a subring of L. It contains A and S by construction and is therefore equal to L. In particular, it contains \(\mathrm{Ke}_1\) , so K is equal to \(\mathrm{A}[a^{-1}]\) .

Now suppose that an element \(s\) of \(S\) is transcendental over \(K\) . Denote by \(E\) the field of fractions of the ring \(A[s]\) . The \(A[s]\) -algebra \(L\) is generated by \(S - \{s\}\) . By the induction hypothesis, there exists a nonzero polynomial \(P \in A[X]\) such that \(E\) is equal to \(A[s][P(s)^{-1}]\) . Let \(\overline{K}\) be an algebraic closure of \(K\) (V, §4, No.~3, p.~23, Theorem 2). Since the field \(\overline{K}\) is infinite (V, §4, No.~1, p.~20, Proposition 3), there exists an element \(x\) of \(\overline{K}\) such that \(P(x) \neq 0\) . Let \(\varphi : E \to \overline{K}\) be the unique homomorphism from the K-algebra \(E = A[s][P(s)^{-1}]\) to the K-algebra \(\overline{K}\) that sends \(s\) to \(x\) . This is absurd because \(E\) is a transcendental extension of \(K\) and \(\overline{K}\) is an algebraic extension of \(K\) . This completes the proof of the theorem.

COROLLARY 1. --- Let K be a commutative field, A be a commutative K-algebra generated by finitely many elements, and m be a maximal ideal of A. a) The degree of A/m over K is finite.

\begin{enumerate}
\def\labelenumi{\alph{enumi})}
\setcounter{enumi}{1}
\tightlist
\item
  Let \(\Omega\) be an algebraically closed extension of K. There exists a K-algebra homomorphism from A to \(\Omega\) with kernel \(\mathfrak{m}\) .
\end{enumerate}

The K-algebra A/m is generated by finitely many elements, and A/m is a field. By Theorem 1, the degree of A/m is finite; assertion a) follows. Every extension of K of finite degree is isomorphic to a subextension of \(\Omega\) (V, §4, No.~1, p.~20, Theorem 1); this gives b).

COROLLARY 2. --- Let K be a commutative field, n be a natural number, \((\mathrm{P}_{i})_{i\in\mathrm{I}}\) be a family of elements of \(K[X_{1},\ldots,X_{n}]\) , and \(\Omega\) be an algebraically closed extension of K. The following conditions are equivalent:

\begin{enumerate}
\def\labelenumi{(\roman{enumi})}
\item
  The polynomials \(\mathrm{P}_i\) do not have a common zero in \(\Omega^n\) .
\item
  There exists a family \((\mathrm{Q}_i)_{i\in \mathrm{I}}\) , with finite support, of elements of \(\mathrm{K}[\mathrm{X}_1,\dots ,\mathrm{X}_n]\) such that \(\sum_{i\in \mathrm{I}}\mathrm{P}_i\mathrm{Q}_i = 1\) .
\end{enumerate}

Denote by A the ring \(K[X_{1},\ldots,X_{n}]\) and by a the ideal generated by the polynomial \(P_{i}\) . Condition (i) means that there are no K-algebra homomorphisms from A/a to \(\Omega\) . If it is satisfied, then by Corollary 1, the ring A/a has no maximal ideals and is therefore zero, and 1 belongs to a. This proves that (i) implies (ii). The implication (ii) \(\Rightarrow\) (i) is clear.

